\documentclass[11pt]{article}
\usepackage[a4paper,margin=27mm]{geometry}
\usepackage{amsmath,amssymb,mathtools}
\usepackage{fontspec,unicode-math}
\setmainfont{Latin Modern Roman}
\setmathfont{Latin Modern Math}
\usepackage{xurl,hyperref}
\hypersetup{colorlinks=true,urlcolor=blue,linkcolor=blue}
\newcommand{\sref}[2]{\hyperlink{source:#1:#2}{[S#2]}}
\newcommand{\cataloguescope}[1]{\paragraph{Mathematical question.} #1}
\begin{document}
\section{Problem 4.6}\label{statement}

\subsection{Definitions and mathematical statement}

\paragraph{Shared notation.}$\mathbb N=\{1,2,\ldots\}$, and $\mathbb Z,\mathbb Q,\mathbb R,\mathbb C$ denote the integers, rationals, reals and complex numbers. Any other conventions are specified locally.


For a nonzero algebraic number $\alpha$, let $f_\alpha(X)=a\prod_{j=1}^d(X-\alpha_j)\in\mathbb Z[X]$ be its primitive irreducible polynomial with positive leading coefficient. Its Mahler measure is $M(\alpha)=a\prod_{j=1}^d\max\{1,|\alpha_j|\}$. For an open $V\subset\mathbb C$, define $L(V)=\inf M(\alpha)^{1/d}$ over nonzero algebraic $\alpha$ which are not roots of unity and whose conjugates all lie outside $V$. Interpret arguments in $(-\pi,\pi]$, and for $0<\theta<\pi$ put $V_\theta=\{z\ne0:|\operatorname{Arg}z|>\theta\}$.
\paragraph{Statement recorded in the supplied source.}
Determine the exact value of the Lehmer--Langevin constant $L(V_\theta)$ as a function of every $\theta\in(0,\pi)$. \sref{statement}{2}
\subsection{Short English statement}
Compute the smallest normalized Mahler measure possible when all conjugates avoid the open sector outside the indicated angle.
\subsection{Sources}
\begin{description}
\item[\hypertarget{source:statement:1}{S1}] ULAM.AI and the UnsolvedMath Contributors, UnsolvedMath: A Curated Collection of Open Mathematics Problems, catalogue entry. CC BY 4.0. \url{https://huggingface.co/datasets/ulamai/UnsolvedMath}.
\item[\hypertarget{source:statement:2}{S2}] Michel Waldschmidt, Open Diophantine Problems, arXiv:math/0312440v2 (24 January 2004), Problem 4.6 and preceding definition of the Lehmer--Langevin constant. \url{https://arxiv.org/abs/math/0312440}.
\end{description}
\label{end:statement}
\end{document}
